\section{Comparison sources and their cancellation}\label{sec:sources}

We now express the errors of the trial potential as sources in the
actual heat equation.  The key fact is a conservation identity: along
an isolated obstacle, the signed source is an endpoint difference of
the profile, up to its cutoff commutator.  Its size is therefore controlled by the obstacle's
displacement, although the total variation of the source counts every
jump.  We establish this identity before using it to estimate either
the scalar or the energy pairing.

Throughout this section we impose the interiority and absence of pins
over $B_{50R_*}(x)$ required in Section~\ref{sec:profiles}.  Constants
are uniform in the finite cube, the allowed pin set, and the injection
slot.  The two traversals and their profiles are those defined there:
$a_s$ is the forward profile and $b_s$ is the reverse profile, each
indexed by its own elapsed time $s\in[0,\beta]$.  Profiles are evaluated
on the full slot lattice before restriction to the current holes.

\subsection{Sources on intervals and at operations}

Fix one of the two directions and write $\varphi_s$ for its profile.
On an open interval containing no instantaneous operation, let
$\mathscr L_s^{\mathrm{act}}$ be the generator on the holes, with the
prescribed killing.  The continuous source is the vector
\begin{equation}\label{eq:continuous-source}
 \mathscr L_s^{\mathrm{act}}(\varphi_s|_{H_s})
             -(\partial_s\varphi_s)|_{H_s}.
\end{equation}
Here $H_s$ denotes the hole set on the slice.  At an instantaneous
transition $K:\R^{H_-}\to\R^{H_+}$, the atomic source, on the outgoing slice, is
\begin{equation}\label{eq:atomic-source}
 K(\varphi_-|_{H_-})-\varphi_+|_{H_+}.
\end{equation}
In reverse time these are the transposed maps, in reverse order.
Thus a reverse atomic source lies before the
operation when viewed in forward time.

Write $s^+$ and $s^-$ for these collections of continuous densities
and atomic vectors in the forward and reverse traversals.  Integrating
against a collection will mean a Lebesgue integral of its continuous
part plus the sum over its atoms.  Every source is supported above
$B_{R+2}(x)$.  There are no pins of $E$ or boundary slots in this region,
so killing creates no additional source there.

There are two types of source.  The first is the cutoff commutator,
defined on the full lattice by
\begin{equation}\label{eq:backgroundsource}
 \gamma_s=\mathscr L\varphi_s-\partial_s\varphi_s,
 \qquad b_*:=\beta^{-1}R^{-3}.
\end{equation}
The uncut profile solves the free heat equation.  Consequently
$\gamma_s$ is supported where $R/3\le r_z\le2R$, and
\begin{equation}\label{eq:backgroundbound}
 |\gamma_s(z,a)|\le Cb_*.
\end{equation}
These assertions hold separately for the smooth part and the point
part of either profile.  Indeed the commutator of $\mathscr L$ with
$\chi$ consists of a cutoff second difference times a profile value,
and products of first differences.  On this annulus,
Lemma~\ref{lem:freebounds} gives values at most $C/(\beta R)$ and
first differences at most $C/(\beta R^2)$, uniformly in elapsed time.
The bounds on the cutoff derivatives prove
\eqref{eq:backgroundbound}.  There is no cutoff atom at the seam,
because the sitewise cutoff commutes with both full-profile seam
averaging and restriction within a site.

The remaining sources come from obstacles.  A removed edge from a
hole $j$ to an obstacle $k$ contributes
\begin{equation}\label{eq:removed-edge-source}
 \tfrac12(\varphi_s(j)-\varphi_s(k))
 \quad\text{at the hole }j.
\end{equation}
If an obstacle moves from $k$ to a hole $j$, the hole moves from $j$
to $k$.  The source on this arriving hole is
\begin{equation}\label{eq:transport-source}
 \varphi_s(j)-\varphi_s(k).
\end{equation}
Exchanges between two obstacles create no source.  At a seam the
source is the difference between the actual average on the remaining
hole slots and the full-profile average, restricted to the outgoing
holes.  It vanishes at sites with no obstacles.

Assign a removed-edge source to its obstacle and a transport source
to the obstacle responsible for the transport.  Assign seam sources
at a site to its obstacle strands, in any fixed manner.  If a strand
is isolated, it is the only obstacle there and receives the entire
seam source at that site.  These assignments can be made either in
a block or in a bin, with strands enumerated at its initial slice
in the direction being used.
On $\mathcal G$, tracing a strand back from a source site in
$B_{R+2}(x)$ uses at most $CL$ steps in its block, by
Proposition~\ref{prop:good}.  Thus every strand contributing a
nonzero source starts in $B_{4R}(x)$, for large $\beta$.
All subsequent strand sums are restricted to this ball.

\subsection{Conservation along an isolated strand}

\begin{lemma}[Signed source identity]\label{lem:signedsource}
Fix a good picture in $\mathcal G$ and one directed block $I=[u,v]$.
Let $Z_s$ be an isolated obstacle strand starting above $B_{4R}(x)$
on this block, with its specified sides at operations.  For either
profile part, the total signed mass of the assigned obstacle sources,
obtained by summing over source slots and integrating in time, is
\begin{equation}\label{eq:signedsource}
 \varphi_v(Z_v)-\varphi_u(Z_u)
                  +\int_u^v\gamma_s(Z_s)\,ds.
\end{equation}
The endpoint values include any seam assigned to the block.  The same
formula holds in reverse elapsed time.
\end{lemma}

\begin{proof}
At every smooth time, all neighboring slots of $Z_s$ are holes, by
isolation.  Summing \eqref{eq:removed-edge-source} over these slots
therefore gives
\[
 \frac12\sum_{j\sim Z_s}
             (\varphi_s(j)-\varphi_s(Z_s))
 = (\mathscr L\varphi_s)(Z_s)
 = (\partial_s\varphi_s)(Z_s)+\gamma_s(Z_s).
\]
The transport contribution \eqref{eq:transport-source} is exactly
the profile increment along each jump of $Z_s$.

It remains to check the seam, where the two hole sets can differ.
At the strand's site let $H_-,H_+$ be the incoming and outgoing holes,
and let $Z_-,Z_+$ be the corresponding obstacle slots.  The actual
hole map preserves the sum over holes.  The full-profile seam
preserves the sum over all slots.  Since the obstacle is the only
one at this site, the total seam source is exactly
\begin{align*}
 \sum_{H_+}(K\varphi_- -\varphi_+)
 &=\sum_{H_-}\varphi_- -\sum_{H_+}\varphi_+\\
 &=\varphi_+(Z_+)-\varphi_-(Z_-).
\end{align*}
Thus it also supplies the increment along the obstacle strand.
Integrating the time derivatives and summing all increments
telescopes to \eqref{eq:signedsource}.  Reverse traversal has the
same conservation of counting sums and uses the inverse transport,
so this proof applies with reverse elapsed time as well.
\end{proof}

We next turn this identity into bounds that retain both the signed
mass and the total variation.  Keeping the two quantities separate
is essential: total variation controls propagation and localization,
whereas the smaller signed mass controls long spatial scales.

\begin{lemma}[Block source bounds]\label{lem:cancellation}
On a picture in $\mathcal G$, let a strand start above
$z\in B_{4R}(x)$ in a directed block of length between $L$ and $2L$.
For the smooth profile part, take $\mathcal W=W$.
For a point-profile part, suppose that $d=\sqrt{1+s}$ is comparable
throughout the block and $\inf d\ge L$, and take $\mathcal W=W_d$
at its starting elapsed time.  Then the obstacle sources assigned
to this strand are supported within distance $CL$ of $z$ and have
total variation at most
\begin{equation}\label{eq:blocksourcevariation}
 CL\mathcal W(z).
\end{equation}
If the strand is isolated, their signed sum has absolute value at most
\begin{equation}\label{eq:blocksourcecancellation}
 CL^{3/4}\mathcal W(z).
\end{equation}
Both assertions also hold with $\mathcal W=W$ for the point profile
on blocks in the second half of elapsed time, and for the sum of that
part and the smooth part, after changing the constant.
\end{lemma}

\begin{proof}
The step bounds in Proposition~\ref{prop:good} keep the strand and
its adjacent source sites within $CL$ of its starting site.  The
weights $W$ are comparable throughout this region, because $L\ll D_*$.
The stated hypotheses give the same comparability for $W_d$.
Each strand has bounded degree, contributes for at most $2L$ units
of continuous time, and makes at most $CL$ jumps.  The difference
bounds of Lemma~\ref{lem:freebounds}, together with
\eqref{eq:removed-edge-source} and \eqref{eq:transport-source},
give \eqref{eq:blocksourcevariation}.  At a seam, fixed slot
multiplicity bounds the number of terms, and the difference of the
two averages is bounded by the slot-difference bound.  This adds at
most $C\mathcal W(z)$.

For an isolated strand use Lemma~\ref{lem:signedsource}.  Its endpoint
displacement is at most $CL^{3/4}$.  The spatial endpoint difference
therefore costs at most $CL^{3/4}W(z)$ in the smooth case.  Time
variation costs $CLW(z)/D_*$, and a seam increment costs at most
$CW(z)$.  On the cutoff annulus,
$b_*\le CW(z)/D_*$, so the integral in
\eqref{eq:signedsource} has the same time bound.  Since
$1+L/D_*\le CL^{3/4}$, this proves
\eqref{eq:blocksourcecancellation} for the smooth part.

For the point part the time derivative is bounded by $CW_d/d$.
On the cutoff annulus its commutator is bounded by
\[
 C\bigl(W_d/R+P_d/R^2\bigr)\le CW_d/d,
\]
because $d\le\sqrt{1+\beta}\ll R$.  The same endpoint argument
gives $C[L^{3/4}+1+L/d]W_d(z)$, which proves the stated bound.
Finally Lemma~\ref{lem:freebounds} bounds the point part by the
smooth envelopes when $s\ge\beta/2$.  Addition proves the last
assertion; the division at $\beta/2$ is a block endpoint.
\end{proof}

Near elapsed time zero, $W_d$ varies too quickly to use block
cancellation.  The bins of length between $1$ and $3$ provide the
following uniform replacement.  We write
$H_\beta=C(1+\log\beta)^C$ for a fixed power of a logarithm;
the constants, and hence this power, may be enlarged at subsequent
uses.

\begin{lemma}[Bin source bounds]\label{lem:sourcebins}
On a picture in $\mathcal G$, the obstacle sources assigned to one
strand in one directed bin, starting above $z\in B_{4R}(x)$ at elapsed time $s$,
have total variation at most
\begin{equation}\label{eq:binsourcevariation}
 H_\beta\bigl(W(z)+W_d(z)\bigr),\qquad d=\sqrt{1+s}.
\end{equation}
The two profile parts separately satisfy this bound with their own
weight.  Their source sites lie within $C\log\beta$ of $z$.
\end{lemma}

\begin{proof}
Proposition~\ref{prop:good} bounds both displacement and the number
of jumps in a bin by $C\log\beta$.  Continuous contributions have
bounded time length.  Within a bin, $1+s$ changes by a bounded
factor, while a spatial shift of $C\log\beta$ changes any of the
polynomial envelopes in Lemma~\ref{lem:freebounds} by at most a
fixed power of $1+\log\beta$.  The formulas for removed edges,
transports, and seam averaging now give the result.
\end{proof}

There are at most $CR^3$ possible starting slots in each bin and
$C\beta$ bins, while $W+W_d\le C$.  Including the background,
the total variation of either complete source collection on
$\mathcal G$ is therefore at most $CH_\beta\beta R^3$.
In particular, multiplying a rapidly small propagation error by
one or two such total variations still gives a rapidly small error.

\subsection{Propagation and the two-source identity}

Let $T_{v,u}$ denote actual forward propagation between specified
slices, including exactly the operations lying between them.
Reverse propagation is the adjoint with the ordered endpoints.
Propagate every source from its insertion slice to its terminal cut
in its own direction.  The resulting signed masses are precisely
the residuals defined in Section~\ref{sec:capacity}:
\begin{equation}\label{eq:residualdefinition}
 \mu_+=\int T_{\beta,s}\,s^+(ds)=Uh-h+c\eta,
 \qquad
 \mu_- =\int T^{\mathrm{rev}}_{\beta,s}\,s^-(ds)
                         =U^*h-h+c\eta.
\end{equation}
For example, on a smooth interval the difference between the actual
mass and the profile satisfies the actual heat equation with forcing
\eqref{eq:continuous-source}.  At an operation it is multiplied by
the actual map and then receives \eqref{eq:atomic-source}.
Telescoping from the initial profile $h$ to the terminal profile
$h-c\eta$ proves \eqref{eq:residualdefinition}.

For a scalar pairing it is useful to make this substitution twice.
Read the reverse profile at forward physical time $t$ by writing
$\bar b_t=b_{\beta-t}$, always on the specified side of an operation.
In particular it is evaluated on the outgoing slice when paired
with a forward source.

\begin{lemma}[Two-source Duhamel identity]\label{lem:duhamel}
For the comparison sources just defined,
\begin{equation}\label{eq:doubledduhamel}
 \ip{\mu_+}{h}
 =\int\ip{s^+(ds)}{\bar b_s}
  +\iint_{s\prec t}
          \ip{T_{t,s}s^+(ds)}{s^-(dt)}.
\end{equation}
In the second integral the reverse sources are indexed by forward
physical time.  The order $s\prec t$ means that the forward insertion
slice lies before the reverse insertion slice.  Thus $T_{t,s}$ runs
from after the first operation to before the later operation when
both sources are atomic.  There is no term pairing the forward and
reverse sources of the same operation.
\end{lemma}

\begin{proof}
Here is the exact telescoping calculation for instantaneous factors;
it also fixes the endpoint convention in the integrals.  Let
$K_j:V_{j-1}\to V_j$, $1\le j\le n$, be the chronological factors
between counting spaces, and write $a_j,b_j$ for comparison values
at those slices, with $a_0=b_n=h$.  In this calculation $b_j$
is indexed in forward order.  Put
\[
 p_j=K_j a_{j-1}-a_j,\qquad
 q_j=K_j^*b_j-b_{j-1},\qquad
 P_{r,l}=K_r\cdots K_{l+1}.
\]
The forward error is $\sum_jP_{n,j}p_j$.  Reverse telescoping gives
\[
 P_{n,j}^*h=b_j+\sum_{k>j}P_{k-1,j}^*q_k.
\]
Their pairing is consequently
\[
 \sum_j\ip{p_j}{b_j}
       +\sum_{j<k}\ip{P_{k-1,j}p_j}{q_k}.
\]
The indices $j<k$ express that $p_j$ is inserted after its factor,
whereas $q_j$ is inserted before that factor.  On each intervening
smooth interval the variation-of-constants formula gives the same
identity with time integrals.  Composing these formulas with the
finite telescoping calculation proves \eqref{eq:doubledduhamel}.
Two continuous source measures have no diagonal mass; atomic terms
retain the strict slice order above.  In particular the forward
terminal seam source pairs directly with $h$ in the first sum.
The reverse seam source lies before that operation and pairs only
with earlier forward sources, through kernels ending before the
seam.  No kernel in this pairing crosses that seam twice.
\end{proof}
